% Generated draft: requires mathematical and source review.
\documentclass[11pt]{article}
\usepackage[margin=25mm]{geometry}
\usepackage{fontspec}
\setmainfont{Latin Modern Roman}
\usepackage{amsmath,amssymb,mathtools,unicode-math}
\setmathfont{Latin Modern Math}
\usepackage{enumitem,needspace,xurl,hyperref}
\setlength{\parindent}{0pt}
\setlength{\parskip}{5pt}
\setlength{\emergencystretch}{3em}
\newcommand{\N}{\mathbb{N}}
\newcommand{\Z}{\mathbb{Z}}
\newcommand{\Q}{\mathbb{Q}}
\newcommand{\R}{\mathbb{R}}
\newcommand{\C}{\mathbb{C}}
\newcommand{\F}{\mathbb{F}}
\newcommand{\A}{\mathbb{A}}
\newcommand{\PP}{\mathbb{P}}
\newcommand{\E}{\mathbb{E}}
\newcommand{\eps}{\varepsilon}
\newcommand{\dd}{\,\mathrm{d}}
\begin{document}
\section*{UnsolvedMath ID: 30000840}
\textbf{Source problem code:} catalogue entry
\subsection{Definitions and mathematical statement}
The source asks for a theory rather than stating a single existence theorem with specified hypotheses. Let $(M,\mathcal F)$ denote a foliated manifold. The requested task is to develop a general index theory for foliations which incorporates index bundles whenever those bundles exist and extends the usual index theory of families.

For comparison, in a standard compact-family setting a family of elliptic operators $D_b:E_b^+\to E_b^-$, parametrized by a space $B$, has a family index in $K^0(B)$. When its kernels and cokernels form vector bundles, this index is represented by
\[
\operatorname{Ind}(D)=[\ker D]-[\operatorname{coker}D].
\]
This describes a conventional meaning of an index bundle, not additional hypotheses supplied for the foliation problem. The requested foliation theory should recover the appropriate family index when the foliation arises from such a family, and should account for index bundles in the foliation cases in which they exist. The record does not specify a class of foliations or operators, conditions ensuring bundle existence, an index target, or the precise compatibility maps or index theorem to be proved. These data must be supplied before this research request becomes a uniquely quantified theorem.

The supplied literature assessment, checked on 2026-08-21, labels the problem partially solved. It attributes index-bundle frameworks for Riemannian foliations to higher harmonic signature and basic index theorems, while saying that a fully general theory for all foliations and all index bundles is not established. These are dated source assessments, not independently verified conclusions of this draft.
\subsection{Short English statement}
Develop an index theory for general foliations that extends family index theory and includes index bundles whenever they exist. The source leaves the admissible operators, bundle conditions and index targets unspecified; the supplied partial-progress assessment concerns Riemannian foliations.
\subsection{Sources}
\begin{itemize}
\item Local UnsolvedMath registry, permanent ID 30000840, source code catalogue entry; supplied complete source record. The record labels the statement exact and reports checking its content against the official report, with normalization of notation and question form.
\item Noncommutative Geometry (2008), Oberwolfach Reports, DOI 10.4171/owr/2007/43, pp. 2543--2608: \url{https://doi.org/10.4171/owr/2007/43}. The registry gives proposed year 2007.
\item Benameur--Heitsch: \url{https://arxiv.org/abs/0711.0352}, supplied as a reference for the dated partial-progress assessment.
\item Foliated basic index theorem: \url{https://arxiv.org/abs/1008.2016}, also supplied for that assessment.
\end{itemize}
\end{document}
